\chapter{Описание стенда и подготовка}
\label{ch:stand}

\section{Программная среда}

Работа выполняется в облачном редакторе \textbf{Figma}, доступном двумя
способами:
\begin{itemize}
  \item \textbf{Веб-версия} --- открыть \texttt{https://www.figma.com}
        в современном браузере (рекомендуется Chromium-семейство:
        Google~Chrome, Microsoft~Edge, Yandex~Browser);
  \item \textbf{Desktop-приложение} --- установить клиент для
        Windows/macOS со страницы загрузки. Функционал идентичен
        веб-версии; единственное преимущество --- работа в нативном окне
        и поддержка локальных шрифтов системы.
\end{itemize}
Для выполнения данной лабораторной работы достаточно браузерной версии.
Регистрация в Figma выполняется по электронной почте или через аккаунт
Google. Бесплатный тарифный план \texttt{Starter} позволяет создавать
до трёх файлов в одной команде, чего достаточно для учебного проекта.

% — Скриншот: общий вид рабочего пространства Figma после входа —
\begin{figure}[h]
  \centering
  \includegraphics[width=0.85\textwidth]{img/01_workspace}
  \caption{Рабочее пространство Figma после входа в аккаунт}
  \label{fig:workspace}
\end{figure}

\section{Подготовка рабочего пространства}

Этапы:
\begin{enumerate}
  \item На главной странице Figma в команде \texttt{Drafts} нажать
        \texttt{New design file}. Откроется новый пустой файл.
  \item Переименовать файл в \texttt{Lab~10~--~Plant~Shop}
        (двойной клик по названию вверху окна).
  \item На левой панели страница по умолчанию называется
        \texttt{Page~1}. Переименовать её в \texttt{Главная страница}.
\end{enumerate}

\section{Создание десктоп-фрейма}

Десктоп-макет создаётся как фрейм фиксированной ширины:
\begin{enumerate}
  \item Выбрать инструмент \texttt{Frame}~(\texttt{F}).
  \item На правой панели в разделе \texttt{Design~\(\rightarrow\)~Desktop}
        выбрать \texttt{Desktop~--~1440\(\times\)1024} или вручную
        нарисовать прямоугольник нужного размера, затем растянуть его
        по высоте.
  \item Задать имя фрейма \texttt{Desktop~1440} (двойной клик по имени
        в левой панели слоёв).
  \item В разделе \texttt{Fill} задать фон \texttt{\#FFFFFF}
        (или светлый зелёный оттенок \texttt{\#F5F8F2}, тематически
        ближе к магазину растений).
\end{enumerate}

% — Скриншот: созданный фрейм Desktop 1440 на холсте —
\begin{figure}[h]
  \centering
  \includegraphics[width=0.85\textwidth]{img/02_frame_1440}
  \caption{Десктоп-фрейм \texttt{Desktop~1440} с белым фоном}
  \label{fig:frame_1440}
\end{figure}

\section{Установка плагина Unsplash}

Плагин \texttt{Unsplash} даёт прямой доступ к каталогу бесплатных
фотографий из Figma.

\begin{enumerate}
  \item Открыть меню \texttt{Main~menu~\(\rightarrow\)~Plugins~\(\rightarrow\)
        Manage plugins}.
  \item Перейти на вкладку \texttt{Browse plugins in Community}.
  \item В строке поиска ввести \texttt{Unsplash}; выбрать одноимённый
        плагин разработчика \textit{Unsplash}.
  \item Нажать \texttt{Install} (потребуется быть авторизованным
        в Figma).
\end{enumerate}

После установки плагин доступен через
\texttt{Resources~(Shift+I)~\(\rightarrow\)~Plugins~\(\rightarrow\)~Unsplash}
или контекстным меню по правому клику на выделенной фигуре:
\texttt{Plugins~\(\rightarrow\)~Unsplash}.

\section{Эталонные примеры}

В тексте методички (\texttt{Методичка~(lab10).pdf}) приведены ссылки на
три эталонных макета:
\begin{itemize}
  \item \textbf{Пример~1} --- общая структура страницы интернет-магазина;
  \item \textbf{Пример~2} --- блок <<Предложения месяца>> и футер;
  \item \textbf{Пример~3} --- адаптивные версии (\(768\,\)px и
        \(360\,\)px).
\end{itemize}
Эталоны открыты в режиме просмотра и используются как ориентир
по типографике, цветовой палитре и расположению блоков. Точное
повторение не требуется: допускается выбор тематики фотографий и
цветовой схемы по вкусу автора.

\endinput
